\chapter{Conclusiones}

El presente Trabajo de Fin de Grado ha abordado el diseño, la implementación y la documentación de un videojuego perteneciente al subgénero \textit{rogue-like}, desarrollado mediante el motor Unity y el lenguaje C\#. El objetivo principal ha sido aplicar de manera práctica principios de ingeniería del software en un contexto real de desarrollo de videojuegos, prestando especial atención a la modularidad del sistema, la generación procedural de contenido y la automatización de procesos dentro del entorno de desarrollo.

Los resultados obtenidos permiten afirmar que se ha logrado construir un sistema software funcional y completo, capaz de soportar un ciclo de juego íntegro desde el menú inicial hasta las condiciones de victoria o derrota. La arquitectura implementada ha favorecido una clara separación de responsabilidades entre subsistemas —gestión de estados, generación de salas, combate, enemigos e interfaz de usuario—, contribuyendo a la mantenibilidad y extensibilidad del proyecto.

La generación procedural de salas ha demostrado ser una solución eficaz para garantizar la rejugabilidad y, al mismo tiempo, un elemento relevante desde el punto de vista de la ingeniería del software, al requerir algoritmos capaces de asegurar coherencia espacial y estabilidad del sistema. Asimismo, la implementación de enemigos con comportamientos diferenciados y la utilización de enemigos especiales basados en escalado de atributos ha permitido introducir progresión de dificultad sin aumentar de forma significativa la complejidad del sistema.

El desarrollo de herramientas específicas integradas en el editor de Unity ha tenido un impacto positivo en la productividad y en la reducción de errores de configuración, evidenciando la importancia del diseño de herramientas auxiliares como parte del proceso de desarrollo de software interactivo.

El proyecto presenta limitaciones derivadas principalmente del alcance temporal y de su desarrollo individual. Entre ellas se incluyen la ausencia de un sistema completo de audio, la falta de persistencia avanzada de progreso entre partidas, la implementación de inteligencia artificial basada en comportamientos simples y la inexistencia de pruebas automatizadas o evaluaciones con usuarios externos. Estas limitaciones no comprometen la validez del sistema desarrollado, pero sí definen líneas claras de trabajo futuro.

En conjunto, este trabajo ha permitido consolidar competencias fundamentales de la Ingeniería Informática, tales como el diseño de arquitecturas software, la implementación de sistemas interactivos en tiempo real, el uso de herramientas profesionales de desarrollo y la documentación técnica de proyectos complejos. Por todo ello, se considera que los objetivos planteados inicialmente han sido alcanzados de forma satisfactoria.

Desde el punto de vista de la práctica profesional en el ámbito de la ingeniería informática, el proyecto evidencia la viabilidad de aplicar metodologías iterativas e incrementales en el desarrollo individual de videojuegos, así como la utilidad de la automatización de tareas dentro del entorno de desarrollo para mejorar la productividad y reducir errores.
